\chapter{Stand der Forschung}
\label{ch:chapter03}
Dieses Kapitel stellt zwei zentrale Arbeiten vor, die für die vorliegende Untersuchung 
von besonderer Relevanz sind. Die erste Arbeit vergleicht den Overhead verschiedener Distributed-Tracing-Tools 
in AWS-Lambda-Funktionen. 
Die zweite Arbeit evaluiert die Performance unterschiedlicher Programmiersprachen auf derselben Plattform. 
Beide Studien werden hinsichtlich ihrer Motivation, Methodik, Ergebnisse und Limitationen analysiert, 
um darauf aufbauend die Forschungslücke für die eigene Arbeit zu adressieren.

\section{Vergleichsstudie 1}
\label{sec:a-comparison}
\emph{\glqq A Comparison of Distributed Tracing Tools in Serverless Applications\grqq }
\subsection{Motivation}
\label{subsec:eder-motivation}

Die Autoren adressieren die wachsende Bedeutung von Observability in serverlosen Anwendungen. 
Da diese aufgrund ihrer feingranularen, verteilten Natur komplex und fehleranfällig sein können, 
ist Distributed Tracing ein zentrales Werkzeug zur Nachverfolgung von Anfragen. 
Gleichzeitig folgt Serverless Computing einem nutzungsbasierten Abrechnungsmodell, 
sodass der durch Tracing verursachte Overhead direkte Auswirkungen auf die Kosten hat. 
Die Autoren motivieren ihre Arbeit mit dem Fehlen eines detaillierten Vergleichs verschiedener Tracing-Tools 
speziell in serverlosen Umgebungen~\cite{10254754}.

\subsection{Untersuchungsgegenstand und Methodik}
\label{subsec:eder-methodology}

Die Studie untersucht drei Open-Source-Distributed-Tracing-Tools – Zipkin, OpenTelemetry (mittels AWS Distro for OpenTelemetry, ADOT) 
und Apache SkyWalking. Alle Tools wurden in AWS-Lambda-Funktionen integriert, wobei Node.js 16.x als konkrete Laufzeitumgebung 
gewählt wurde. Die Evaluation erfolgte anhand von fünf Microbenchmark, die typische Interaktionsmuster serverloser Anwendungen abbilden 
(siehe Unterabschnitt~\ref{subsec:lambda-benchmarks}).

Für jede Kombination aus Benchmark und Tracing-Tool wurden 100 warme Funktionsaufrufe 
(hot starts) durchgeführt und die Metriken Laufzeit (Runtime) und Speicherverbrauch (Max Memory Used) erfasst.
Die Funktionsaufrufe wurden in 10 Batches zu je 10 Aufrufen organisiert und nach jedem Batch wurden die Lambda-Funktionen redeployed. 
Zusätzlich wurde eine nicht-instrumentierte Baseline gemessen.
Für die Messung des Overheads bei kalten Funktionsaufrufen (cold starts) wurde eine leere Funktion verwendet. 
Diese leere Funktion wurde in 4 Varianten implementiert - ohne Tracing, mit Zipkin, mit OpenTelemetry und mit SkyWalking. 
Jede Variante wurde 100 Mal aufgerufen (100 Kaltstarts), um die Metrik Init Duration zu messen.
Die Tracing-Backends wurden auf einer AWS-EC2-Instanz gehostet.

Für alle Metriken wurden statistische Tests (einseitige Welch's t-Tests) durchgeführt, um zu bestimmen, ob die Unterschiede 
zwischen den Tracing-Tools und der Baseline statistisch signifikant sind.
\subsection{Ergebnisse}
\label{subsec:overhead-results}


\definecolor{tablegray}{gray}{0.9}

\begin{table}[H]
\centering
\caption{Effizienz-Overhead Vergleich zwischen Zipkin, OTel und SkyWalking. Resultate entnommen aus~\cite{10254754}.}
\label{efficiency-overhead}
\begin{tabular}{llccc}
\toprule
 & & \textbf{Zipkin} & \textbf{OTel} & \textbf{SkyWalking} \\
\midrule
\multirow{5}{*}{Runtime} & Function invocation & 6.37 \% & 4.52 \% & \cellcolor{tablegray}38.29 \% \\
\cmidrule{2-5}
 & HTTP GET & \cellcolor{tablegray}36.08 \% & \cellcolor{tablegray}70.85 \% & \cellcolor{tablegray}48.09 \% \\
 & HTTP POST & 6.31 \% & 17.15 \% & 8.15 \% \\
\cmidrule{2-5}
 & Database read & 4.88 \% & 7.17 \% & \cellcolor{tablegray}73.04 \% \\
 & Database write & 0 \% & 21.25 \% & \cellcolor{tablegray}85.79 \% \\
\midrule
Memory Usage & & \cellcolor{tablegray}7.51 \% & \cellcolor{tablegray}91.06 \% & \cellcolor{tablegray}18.90 \% \\
\midrule
Initialization & & \cellcolor{tablegray}70.22 \% & \cellcolor{tablegray}575.35 \% & \cellcolor{tablegray}208.04 \% \\
duration & & & & \\
\bottomrule
\end{tabular}
\end{table}

Tabelle~\ref{efficiency-overhead} zeigt die Effizienz-Overheads der Tracing-Tools. Die graumarkierten Zellen heben statistisch signifikante Overheads hervor.

Zipkin verursachte den geringsten durchschnittlichen Laufzeit-Overhead (10,73\%), SkyWalking den höchsten (50,67\%), OTel lag bei 24,19\%. SkyWalking war das einzige Tool mit signifikantem Laufzeit-Overhead über alle Benchmarks hinweg, allerdings erreichte OTel bei einem Benchmark höhere Werte.

Für Speicherverbrauch und Initialisierungszeiten zeigten sich über alle Benchmarks signifikante Steigerungen. Der Speicherverbrauch reichte von 7,51\% (Zipkin) bis 91,06\% (OTel), die Initialisierungsdauer von 70,22\% (Zipkin) bis 575,35\% (OTel). 

Insgesamt erwies sich der frameworkbasierte Ansatz von Zipkin als am effizientesten, während die agentenbasierten Ansätze von OTel und SkyWalking deutlich höhere Overheads verursachten.

\subsection{Einordnung}
\label{subsec:eder-discussion}

Die Studie von Eder et al. liefert erstmals einen systematischen Vergleich von Tracing-Tools in einer serverlosen 
Umgebung. Die Ergebnisse zeigen, dass der Overhead erheblich variieren kann. Besonders auffällig ist der hohe 
Speicher- und Kaltstart-Overhead von ADOT, der von den Autoren auf den zusätzlichen Collector zurückgeführt wird.
Allerdings beschränkt sich die Untersuchung auf Node.js als Laufzeitumgebung. 
Zudem wurden bei den Kaltstarts leere Funktionen verwendet, was den 5 Microbenchmark nicht entspricht. 
Die Autoren selbst regen an, weitere Sprachen und Plattformen in zukünftige Vergleiche einzubeziehen – eine 
zentrale Motivation für die vorliegende Arbeit.

\section{Vergleichsstudie 2}
\label{sec:rodriguez-2023}
\emph{\glqq Evaluation of Programming Languages for Memory Usage, Scalability, and Cold Start, on AWS Lambda Serverless Platform as a Case Study\grqq }

\subsection{Motivation}
\label{subsec:rodriguez-motivation}

Diese Studie von Rodriguez et al.~\cite{Rodriguez2024Evaluation} adressiert die Frage, wie
 sich verschiedene Programmiersprachen in serverlosen Umgebungen 
 hinsichtlich Initialisierung, Kaltstarts, Speichernutzung und Skalierbarkeit verhalten. 
 Die Autoren stellen fest, dass viele wissenschaftliche Arbeiten die Analyse dieser Faktoren in Bezug auf die Sprachwahl vernachlässigen. 
 Ziel der Arbeit ist es, eine Entscheidungsgrundlage für Entwickler zu schaffen, 
 die serverlose Anwendungen auf AWS Lambda implementieren.

\subsection{Untersuchungsgegenstand und Methodik}
\label{subsec:rodriguez-methodology}

Die Studie vergleicht drei Programmiersprachen: JavaScript (Node.js), Python und Java. 
Als Anwendungsfall dient eine REST-API für Create-, Read-, Update- und Delete-Operationen auf DynamoDB. 
Die API wurde in allen drei Sprachen implementiert und mit dem Open-Source-Tool Artillery getestet. 

Die Experimente umfassen Initialisierungstests, bei denen fünf Kaltstarts pro Sprache und CRUD-Operation bei 10 GB 
Speicherzuweisung gemessen wurden. Zusätzlich wurde im Anschluss an jeden Kaltstart jeweils eine warme Ausführung gemessen, 
um den Unterschied zwischen Kalt- und Warmstart zu quantifizieren. Weiterhin wurden Skalierbarkeitstests durchgeführt, 
bei denen Anfragen mit steigender Anzahl pro Sekunde (1, 10, 25, 50, 100, 200 gleichzeitige Aufrufe) 
zur Beobachtung des Skalierungsverhaltens der Leseoperation (Read) eingesetzt wurden. 
Abschließend erfolgten Speichertests durch Variation des zugewiesenen Speichers (512 MB, 1, 2, 4, 6, 8, 10 GB) 
zur Analyse des Einflusses auf Laufzeit und Initialisierung, beispielhaft für die Create-Operation.

\subsection{Ergebnisse}
\label{subsec:rodriguez-results}
\footnotetext{Alle Ergebnisse aus den Tabellen~\ref{tab:rodriguez-init} bis~\ref{tab:rodriguez-scalability} wurden aus~\cite{Rodriguez2024Evaluation} entnommen.}

\paragraph{Initialisierungszeiten}
Tabelle~\ref{tab:rodriguez-init} zeigt die durchschnittlichen Initialisierungszeiten pro CRUD-Operation. Java weist mit Abstand 
die längsten Zeiten auf (bis zu 1630 ms), während Node.js und Python ähnliche Werte im Bereich von 360–440 ms liefern.

\begin{table}[H]
\centering
\caption{Durchschnittliche Initialisierungszeiten pro CRUD-Operation und Programmiersprache (in ms)~\cite{Rodriguez2024Evaluation}}
\label{tab:rodriguez-init}
\begin{tabular}{lcccc}
\hline
\textbf{Sprache} & \textbf{Create} & \textbf{Read} & \textbf{Update} & \textbf{Delete} \\
\hline
Node.js & 418,23 & 431,04 & 413,58 & 442,15 \\
Python  & 384,90 & 400,43 & 361,62 & 411,20 \\
Java    & 1623,53 & 1630,52 & 1526,12 & 1577,30 \\
\hline
\end{tabular}
\end{table}

\paragraph{Laufzeiten (Kalt vs. Warm)}
Die Tabellen~\ref{tab:rodriguez-nodejs} bis~\ref{tab:rodriguez-java} zeigen die durchschnittlichen Ausführungszeiten für 
Kalt- und Warmstarts pro Sprache und Operation. Java verbessert sich am stärksten bei warmen Ausführungen – 
die Laufzeit sinkt um durchschnittlich 84\,\%, während Node.js und Python Verbesserungen um etwa 45–48\,\% erzielen.

\begin{table}[H]
\centering
\caption{Ausführungszeiten in Node.js für Kalt- und Warmstarts (in ms)~\cite{Rodriguez2024Evaluation}}
\label{tab:rodriguez-nodejs}
\begin{tabular}{lcccc}
\hline
\textbf{Status} & \textbf{Create} & \textbf{Read} & \textbf{Update} & \textbf{Delete} \\
\hline
Kalt  & 76,69 & 74,59 & 84,04 & 77,76 \\
Warm  & 51,23 & 39,24 & 45,53 & 41,06 \\
Verbesserung & 33,2\,\% & 47,4\,\% & 45,8\,\% & 47,2\,\% \\
\hline
\end{tabular}
\end{table}

\begin{table}[H]
\centering
\caption{Ausführungszeiten in Python für Kalt- und Warmstarts (in ms)~\cite{Rodriguez2024Evaluation}}
\label{tab:rodriguez-python}
\begin{tabular}{lcccc}
\hline
\textbf{Status} & \textbf{Create} & \textbf{Read} & \textbf{Update} & \textbf{Delete} \\
\hline
Kalt  & 60,73 & 49,75 & 65,73 & 43,97 \\
Warm  & 32,58 & 27,38 & 38,47 & 22,56 \\
Verbesserung & 46,3\,\% & 45,0\,\% & 41,5\,\% & 48,7\,\% \\
\hline
\end{tabular}
\end{table}

\begin{table}[H]
\centering
\caption{Ausführungszeiten in Java für Kalt- und Warmstarts (in ms)~\cite{Rodriguez2024Evaluation}}
\label{tab:rodriguez-java}
\begin{tabular}{lcccc}
\hline
\textbf{Status} & \textbf{Create} & \textbf{Read} & \textbf{Update} & \textbf{Delete} \\
\hline
Kalt  & 701,19 & 664,63 & 644,33 & 581,49 \\
Warm  & 116,64 & 98,10 & 116,89 & 84,65 \\
Verbesserung & 83,4\,\% & 85,2\,\% & 81,9\,\% & 85,4\,\% \\
\hline
\end{tabular}
\end{table}

\paragraph{Skalierbarkeit}
Tabelle~\ref{tab:rodriguez-scalability} zeigt die Laufzeiten der Leseoperation (Read) bei steigender Anzahl paralleler Aufrufe.
Java startet mit einer hohen Latenz (52 ms) aufgrund von Kaltstarts, verbessert sich aber bei 10 Aufrufen  und 
unterbietet damit Node.js und Python. Bei 25 Aufrufen treten in Java erneut Kaltstarts auf, 
während Node.js und Python stabil bleiben. Ab 50 Aufrufen stabilisieren sich alle Sprachen. Bei 200 Aufrufen erreicht 
den besten Wert, gefolgt von Python und Node.js.

\begin{table}[H]
\centering
\caption{Laufzeiten der Leseoperation (Read) bei steigender Anzahl der Anfragen pro Sekunde (in ms)~\cite{Rodriguez2024Evaluation}}
\label{tab:rodriguez-scalability}
\begin{tabular}{lcccccc}
\hline
\textbf{Sprache - Parallelität} & \textbf{1} & \textbf{10} & \textbf{25} & \textbf{50} & \textbf{100} & \textbf{200} \\
\hline
Node.js & 70,3 & 30,3 & 39,9 & 29,7 & 26,2 & 26,6 \\
Python  & 58,1 & 31,6 & 16,0 & 16,5 & 15,0 & 15,4 \\
Java    & 52,7 & 20,0 & 158,3 & 24,3 & 19,7 & 13,9 \\
\hline
\end{tabular}
\end{table}

\paragraph{Einfluss des Speichers auf die Laufzeit}

Geringerer Speicher führt bei allen Sprachen zu längeren Laufzeiten, 
jedoch ist der Effekt bei Java besonders ausgeprägt. Unter 1024 MB steigen die Laufzeiten drastisch an, während Node.js und Python auch 
bei geringerem Speicher akzeptable Werte liefern. Bei 4096 MB gleichen sich die Laufzeiten aller Sprachen an (für mehr Details siehe~\cite{Rodriguez2024Evaluation}).

\subsection{Einordnung}
\label{subsec:rodriguez-discussion}
Die Arbeit bestätigt, dass verschiedene Programmiersprachen unterschiedliche Stärken und Schwächen im Bezug auf 
die Performance aufweisen. Die Studie unterstreicht 
die Notwendigkeit, die Sprachwahl anhand des erwarteten Lastprofils und der verfügbaren Ressourcen zu treffen.

Allerdings sind bei der Interpretation der Ergebnisse einige methodische Einschränkungen und Mängel zu beachten. 
Erstens basieren die Messungen auf lediglich fünf Kalt- und fünf Warmstarts pro Konfiguration, was aus statistischer Sicht eine zu 
geringe Stichprobengröße darstellt und die Aussagekraft der Messwerte erheblich reduziert. 
Zweitens wird nicht dokumentiert, ob und welche reservierte Konkurrenz (reserved concurrency) für die Funktionen konfiguriert war 
– ein Faktor, der das Skalierungsverhalten insbesondere bei parallelen Aufrufen erheblich beeinflussen kann. 
Drittens berücksichtigt die gewählte Methodik nicht die Optimierungsmechanismen der Laufzeitumgebungen. 
Sowohl die \ac{JIT}-Compiler von Java und Node.js als auch der CPython-Interpreter benötigen mehrere Durchläufe, 
um ihren Code zu optimieren~\cite{CPythonPEP659,FromWarmToHotStarts}. Mit nur einem einzigen Warmstart pro Funktion werden diese Effekte nicht erfasst, 
was zu irreführenden Ergebnissen führt. Dies zeigt sich beispielsweise im Vergleich der initialen Warmstart-Messungen mit 
den Skalierbarkeitstests. Während Java in den ersten warmen Aufrufen noch relativ hohe Laufzeiten aufweist, 
erreicht es bei 200 Anfragen pro Sekunde die besten Werte – ein Hinweis darauf, dass die \ac{JIT}-Optimierung erst nach mehreren 
Aufrufen ihre volle Wirkung entfaltet.

Trotz dieser Einschränkungen liefert die Studie wertvolle Hinweise auf die relativen Stärken und Schwächen der untersuchten Sprachen. 
Die Studie \emph{\glqq Cold Start Influencing Factors in Function as a Service\grqq } von Manner et al.~\cite{FaaSColdStart2018} untersucht den 
Kaltstart‑Overhead auf AWS Lambda und Azure Functions für Java und JavaScript. Ihre Ergebnisse bestätigen die hier beobachteten Tendenzen. 
Java verursacht einen deutlich höheren Kaltstart‑Overhead als interpretierte Sprachen (Faktor 2–3) und mehr Speicher reduziert den 
Overhead messbar. Damit bestätigt diese Studie die grundlegenden Zusammenhänge zwischen Sprachwahl, Ressourcenkonfiguration und 
Kaltstartverhalten, die auch in der Arbeit von Rodriguez et al. prägend sind.
Die genannten methodischen Mängel unterstreichen jedoch die Notwendigkeit weiterführender Untersuchungen mit größeren Stichproben 
und realistischeren Ausführungsprofilen.

\subsection{Limitationen und Forschungslücke}
\label{sec:limitations}

Beide vorgestellten Studien liefern wertvolle Erkenntnisse für die serverlose Entwicklung, weisen jedoch spezifische Einschränkungen auf, 
die die Notwendigkeit weiterer Forschung begründen.

\paragraph{Eingeschränkte Untersuchungsvielfalt}
Die Untersuchung von Eder et al.~\cite{10254754} beschränkt sich auf Node.js, während Rodriguez et al.~\cite{Rodriguez2024Evaluation} 
immerhin drei Sprachen vergleichen, aber ohne Berücksichtigung von Tracing-Overhead. Dabei fehlt die Kombination beider Aspekte.

\paragraph{Statistische Aussagekraft}
Die Stichprobengrößen sind teils gering (100 Aufrufe pro Benchmark bei Eder et al. und 5 bei Aufrufe bei Rodriguez et al.), 
was die Verallgemeinerbarkeit einschränkt und die statistische Aussagekraft der Ergebnisse beeinträchtigt.

\section{Fazit}
\label{sec:conclusion-related}

Die betrachteten Vergleichsstudien belegen, dass sowohl die Wahl der Programmiersprache als auch der Einsatz von Distributed-Tracing-Tools 
erheblichen Einfluss auf die Performance serverloser Anwendungen haben. Dabei sind alle 3 Performance-Metriken (Laufzeit, Speicherverbrauch und Initialisierungszeit) betroffen. 
Die Ergebnisse unterstreichen die Notwendigkeit, diese Faktoren bei der Entwicklung und Wartung von serverlosen Anwendungen sorgfältig abzuwägen. 
Jedoch ist dies ohne eine umfassende Analyse der Wechselwirkungen zwischen diesen Faktoren nicht möglich.